\chapter{Fuentes de datos y procedencia}\label{cap:datos}
\begin{sloppypar}

% -------------------------------------------------------------------------------------------------
\section{Propósito}\label{sec:cap13-proposito}

Una valuación vale lo que valen sus datos. Este capítulo documenta de dónde sale cada dato que entra al desk,
en qué unidad llega y en cuál se guarda, qué parte es real y qué parte es inventada para demostración, y
cómo se hace visible esa procedencia: las fijaciones diarias de TIIE-F, SOFR y USDMXN (tabla
\texttt{fixings}), los \emph{snapshots} de mercado (tabla \texttt{snapshots}: el sintético \texttt{MX-DEMO} y el
mixto \texttt{MKT-LIVE}), la historia de factores del riesgo y el respaldo de la base. Los capítulos
\ref{cap:curvas}, \ref{cap:swaps-ois} y \ref{cap:ciclo-vida} usan esos datos; la pantalla de Metodología los
declara y aquí se contrasta esa declaración con el código.

% -------------------------------------------------------------------------------------------------
\section{Notación}\label{sec:cap13-notacion}

Los símbolos comunes están en el apéndice~\ref{ap:simbolos}. En este capítulo:
$y_\theta$ es el rendimiento par del Tesoro de EE.\,UU. al plazo $\theta\in\{\text{1M},\dots,\text{10Y}\}$ (decimal);
$D_\theta$ son los días del pilar (30, 91, 182, 365, \dots); $q_\theta$ es la cotización de la curva SOFR que
resulta (decimal, base ACT/360); $u_i$ es la $i$-ésima observación válida de la serie DEXMXUS de FRED (pesos por
dólar); $n$ es el número de cambios diarios con que se mide una volatilidad; $\hat\sigma$ es una desviación
estándar muestral; $R$ es la matriz de correlaciones de los tres factores, $L$ su factor de Cholesky
($LL^{\top}=R$) y $z_{d,k}$ las normales estándar independientes del día $d$.

% -------------------------------------------------------------------------------------------------
\section{Modelo y fórmulas}\label{sec:cap13-modelo}

\subsection{Unidades}\label{sec:cap13-unidades}

Toda tasa se guarda en decimal ($0.0395=3.95\,\%$). Las fuentes publican en por ciento; la conversión se hace
al leer, en el módulo que parsea cada fuente:
\begin{align}
  r_{\text{SOFR}} &= \operatorname{round}\!\bigl(\texttt{percentRate}/100,\;10\bigr)
    && \text{(Fed de Nueva York)}, \label{eq:cap13-nyfed}\\
  r_{\text{TIIE-F}} &= \operatorname{round}\!\bigl(\texttt{dato}/100,\;10\bigr)
    && \text{(Banxico, SF331451)}, \label{eq:cap13-banxico}\\
  \text{FIX} &= \texttt{dato}
    && \text{(Banxico, SF43718)}, \label{eq:cap13-fix}\\
  y_\theta &= \texttt{rate}/100
    && \text{(MarketDataFeed)}. \label{eq:cap13-mdf}
\end{align}
El FIX está en pesos por dólar y no se convierte. Las ecuaciones~\eqref{eq:cap13-nyfed}--\eqref{eq:cap13-fix} son \codigo{desk/api/nyfed.py}{parse_sofr} y
\codigo{desk/api/banxico.py}{parse_series}; la~\eqref{eq:cap13-mdf} es
\codigo{desk/api/marketdata.py}{latest_treasury_curve} (y \codigo{desk/api/marketdata.py}{latest_sofr}). El
redondeo a 10 decimales es el \texttt{round} de Python. La serie DEXMXUS no se convierte:
\codigo{desk/api/fred.py}{parse_dexmxus} conserva las observaciones con valor finito y positivo y descarta las
vacías (\texttt{""}, \texttt{"."}).

\subsection{Curva SOFR de \texttt{MKT-LIVE}}\label{sec:cap13-curva}

La curva SOFR de \texttt{MKT-LIVE} no viene de swaps: es una \emph{aproximación} con rendimientos par del
Tesoro \citep{treasury2026} puestos en la convención de las cotizaciones OIS de la curva. El pilar de un día
es el SOFR publicado tal cual \citep{nyfed2026}; los demás pilares (1M, 3M, 6M, 1Y, 2Y, 3Y, 5Y, 7Y y 10Y, con
$D_\theta=30,\,91,\,182,\,365,\,730,\,1095,\,1825,\,2555,\,3650$) son
\begin{equation}\label{eq:cap13-tesoro}
  q_\theta=\begin{cases}
    y_\theta\cdot\dfrac{360}{365} & D_\theta\le 365,\\[1.2ex]
    \Bigl[\bigl(1+\tfrac{y_\theta}{2}\bigr)^{2}-1\Bigr]\cdot\dfrac{360}{365} & D_\theta>365,
  \end{cases}
\end{equation}
redondeada a 6 decimales (\codigo{desk/api/live_market.py}{treasury_to_ois},
\codigo{desk/api/live_market.py}{SOFR_PILLARS}). La primera rama traduce una tasa en base 365 a base 360 sin
capitalizar; la segunda trata $y_\theta$ como tasa nominal con capitalización semestral, la lleva a tasa
efectiva anual y luego a base 360. Las cotizaciones resultantes entran al mismo bootstrap OIS que las de
\texttt{MX-DEMO} (capítulo~\ref{cap:curvas}), con depósito el pilar de un día y OIS par el resto.

\begin{diferencia}
  Una curva SOFR construida con swaps OIS descuenta con la tasa de swap; esta usa Treasuries y por tanto omite
  el diferencial entre swaps SOFR y Treasuries, que la Metodología cita como negativo (por ejemplo $-25$~\pb\ a
  2 años y $-87$~\pb\ a 30 años en 2025, según las fuentes que enlaza) y no se mide aquí. La curva de
  \texttt{MKT-LIVE} probablemente sobreestima las tasas de swap en esas magnitudes.
\end{diferencia}

\subsection{Volatilidades diarias de factores}\label{sec:cap13-vols}

El snapshot guarda la volatilidad \emph{diaria} de tres factores (\texttt{TIIE-F}, \texttt{SOFR},
\texttt{USDMXN}) y su matriz de correlaciones. En \texttt{MKT-LIVE} se miden dos de ellas con muestra de a lo
más $n_{\max}=250$ cambios diarios y, como mínimo, $n_{\min}=60$:
\begin{align}
  \hat\sigma_{\text{SOFR}} &= \operatorname{sd}\Bigl\{\,10^{4}\bigl(y^{5\mathrm{Y}}_{i}-y^{5\mathrm{Y}}_{i-1}\bigr)\Bigr\}_{i=1}^{n}
    \quad\text{(pb por día)}, \label{eq:cap13-vol-sofr}\\
  \hat\sigma_{\text{USDMXN}} &= \operatorname{sd}\Bigl\{\,\ln\frac{u_i}{u_{i-1}}\Bigr\}_{i=1}^{n}
    \quad\text{(rendimiento logarítmico por día)}, \label{eq:cap13-vol-fx}
\end{align}
donde $\operatorname{sd}$ es la desviación estándar muestral (divisor $n-1$,
\codigo{desk/api/live_market.py}{sample_std}), la serie de SOFR es el rendimiento par de 5 años del Tesoro
(\emph{no} SOFR) y la de USDMXN une observaciones válidas consecutivas aunque entre ellas haya días sin dato
(\codigo{desk/api/live_market.py}{daily_vol_logret}). La ventana toma las últimas $n_{\max}+1$ observaciones y
con menos de $n_{\min}$ cambios responde 409 y no construye el snapshot
(\codigo{desk/api/live_market.py}{window}, constantes \codigo{desk/api/marketdata.py}{MAX_CHANGES} y
\codigo{desk/api/marketdata.py}{MIN_CHANGES}). Son volatilidades diarias, \emph{no} anualizadas. La
volatilidad del factor \texttt{TIIE-F}, la matriz de correlaciones, la superficie de volatilidad y los
parámetros de Hull--White de \texttt{MKT-LIVE} se heredan sin cambio de \texttt{MX-DEMO} (sintéticos).

\subsection{Procedencia de un snapshot}\label{sec:cap13-procedencia}

Cada snapshot guarda en su JSON un bloque \texttt{sources} con siete componentes (\texttt{curve:TIIE-F},
\texttt{curve:SOFR}, \texttt{spot}, \texttt{surface}, \texttt{hw}, \texttt{factor\_vols},
\texttt{factor\_corr}), cada uno con \texttt{kind}, \texttt{label} y \texttt{as\_of}. El tipo es
\texttt{live}, \texttt{proxy}, \texttt{synthetic} o \texttt{mixed}, y la calidad del snapshot es
\begin{equation}\label{eq:cap13-calidad}
  Q=\begin{cases}
    \text{synthetic} & \text{si todos los componentes son synthetic},\\
    \text{live} & \text{si ninguno es synthetic ni mixed},\\
    \text{mixed} & \text{en otro caso},
  \end{cases}
\end{equation}
(\codigo{desk/api/market.py}{data_quality}); un snapshot sin bloque \texttt{sources} es todo sintético
(\codigo{desk/api/market.py}{sources_of}). La interfaz lo muestra como «Datos sintéticos», «Datos mixtos» o «Datos
de mercado» (\codigo{desk/ui/src/lib/provenance.ts}{qualityText}) y detalla cada componente en un
\emph{tooltip} (\codigo{desk/ui/src/lib/provenance.ts}{sourcesTooltip}). Un componente \texttt{proxy} cuenta como no
sintético en~\eqref{eq:cap13-calidad}, así que un snapshot sólo con componentes \texttt{live} y \texttt{proxy}
se rotularía «Datos de mercado».

\subsection{Datos sintéticos de \texttt{MX-DEMO}}\label{sec:cap13-sinteticos}

Todo \texttt{MX-DEMO} es inventado (\codigo{desk/api/seed_market.py}{demo_data}), con
\texttt{as\_of} $=2026\text{-}11\text{-}10$ (\codigo{desk/api/seed_market.py}{AS_OF}) y
$\Spot=18.60$ (\codigo{desk/api/seed_market.py}{SPOT}). La tabla~\ref{tab:cap13-cotizaciones} lista las
cotizaciones de las dos curvas (la primera de cada una es un depósito a un día y el resto OIS par).

\begin{table}[ht]
  \centering\small
  \caption{Cotizaciones sintéticas de \texttt{MX-DEMO}: TIIE-F (\texttt{\_TIIE\_F}) y SOFR (\texttt{\_SOFR}).}\label{tab:cap13-cotizaciones}
  \begin{tabular}{@{}lrrlrr@{}}
    \toprule
    TIIE-F & días & tasa (\%) & SOFR & días & tasa (\%) \\
    \midrule
    1D & 1 & 7.00 & 1D & 1 & 3.90 \\
    1x28 & 28 & 6.98 & 1M & 30 & 3.85 \\
    3x28 & 84 & 6.92 & 3M & 91 & 3.80 \\
    6x28 & 168 & 6.85 & 6M & 182 & 3.75 \\
    13x28 & 364 & 6.75 & 1Y & 365 & 3.70 \\
    26x28 & 728 & 6.70 & 2Y & 730 & 3.55 \\
    39x28 & 1092 & 6.72 & 3Y & 1095 & 3.50 \\
    65x28 & 1820 & 6.80 & 5Y & 1825 & 3.55 \\
    91x28 & 2548 & 6.90 & 7Y & 2555 & 3.65 \\
    130x28 & 3640 & 7.05 & 10Y & 3650 & 3.80 \\
    \bottomrule
  \end{tabular}
\end{table}

La superficie de volatilidad USDMXN es una función cerrada de la sonrisa: para el vencimiento $j$
(30, 91, 182 y 365 días, con volatilidad al dinero $0.115,\,0.120,\,0.125,\,0.130$) y el precio de ejercicio
$K_i=17.0+0.5\,i$, $i=0,\dots,7$,
\begin{equation}\label{eq:cap13-sonrisa}
  \vol_{j,i}=\vol^{\mathrm{atm}}_{j}+\tfrac12\,k_i^{2}+0.03\,k_i,\qquad k_i=\ln\frac{K_i}{\Spot},
\end{equation}
(\codigo{desk/api/seed_market.py}{_smile}). Los parámetros de Hull--White son $\hwa=0.05$, $\hwsig=0.01$ para
TIIE-F y $\hwa=0.03$, $\hwsig=0.008$ para SOFR; las volatilidades diarias de los factores son 4~\pb\ (TIIE-F),
5~\pb\ (SOFR) y 0.0075 (USDMXN), con correlaciones $\rho_{\text{TIIE,SOFR}}=0.35$,
$\rho_{\text{TIIE,USDMXN}}=0.30$ y $\rho_{\text{SOFR,USDMXN}}=-0.10$. Las cotizaciones de colateral USD
(\codigo{desk/api/seed_market.py}{csa_quotes_data}) son puntos forward de 28, 84, 168 y 364 días
($0.0438$, $0.1317$, $0.2637$ y $0.5745$ pesos por dólar) y una base de cross currency swaps a 2, 3, 5, 7 y
10 años de $-10$, $-12$, $-15$, $-17$ y $-20$~\pb; no hay fuente libre para ninguna de las dos y \texttt{MKT-LIVE}
no las trae. El portafolio de demostración son siete operaciones con contrapartes inventadas (sufijo
«sintético», \codigo{desk/api/seed_trades.py}{TRADES}); el corte de XVA añade otras
(capítulo~\ref{cap:xva}).

La historia sintética de factores de \texttt{MX-DEMO} (\codigo{desk/api/risk/factor_history.py}{seed_synthetic})
son 500 días hábiles (lunes a viernes) anteriores al \texttt{as\_of}, con la semilla fija $20261110$
(\codigo{desk/api/risk/factor_history.py}{SYNTHETIC_SEED}); para cada día $d$ y cada factor $f$
(en el orden \texttt{TIIE-F}, \texttt{SOFR}, \texttt{USDMXN}),
\begin{equation}\label{eq:cap13-historia}
  v_{f,d}=\hat\sigma_f\sum_{k\le f}L_{fk}\,z_{d,k},
\end{equation}
con $\hat\sigma_f$ la volatilidad diaria del snapshot (pb para las tasas, rendimiento logarítmico para
USDMXN) y $z_{d,k}$ de \texttt{random.Random(20261110).gauss(0, 1)}: las mismas cifras en cualquier base.

\subsection{Fijaciones: no se sobrescribe}\label{sec:cap13-fijaciones}

La tabla \texttt{fixings} guarda $(\text{índice},\text{fecha},\text{valor})$ con clave primaria
$(\text{índice},\text{fecha})$; las tasas van en decimal y \texttt{USDMXN} es el FIX en pesos por dólar
(\codigo{desk/api/fixings.py}{INDICES}). Una fijación ya guardada no se cambia nunca: al traer una de la
fuente se compara con la guardada y es
\begin{equation}\label{eq:cap13-conflicto}
  \text{sin cambio si }\;\bigl|v_{\text{fuente}}-v_{\text{guardado}}\bigr|\le 10^{-9},\qquad
  \text{conflicto en otro caso},
\end{equation}
y el conflicto se lista (\texttt{date}, \texttt{stored}, \texttt{source}) sin tocar el valor guardado
(\codigo{desk/api/fixings_fetch.py}{TOLERANCE}). Las fijaciones se usan en el devengo del periodo en curso de
un swap (capítulos~\ref{cap:swaps-ois} y~\ref{cap:ciclo-vida}), en la conversión a pesos de los flujos
liquidados con el FIX de su fecha (capítulo~\ref{cap:ciclo-vida}) y en los avisos y mediciones de las coberturas
(capítulo~\ref{cap:coberturas}). Los índices tienen un calendario de días inhábiles
(capítulo~\ref{cap:introduccion}) que decide si la falta de una fijación es un dato faltante.

% -------------------------------------------------------------------------------------------------
\section{Implementación}\label{sec:cap13-implementacion}

\subsection{Fuentes y módulos}\label{sec:cap13-fuentes}

\begin{table}[ht]
  \centering\small
  \caption{Fuentes externas de datos.}\label{tab:cap13-fuentes}
  \begin{tabular}{@{}>{\raggedright\arraybackslash}p{2.5cm}>{\raggedright\arraybackslash}p{3.5cm}>{\raggedright\arraybackslash}p{2.9cm}>{\raggedright\arraybackslash}p{1.8cm}>{\raggedright\arraybackslash}p{2.8cm}@{}}
    \toprule
    Fuente & Dato & Unidad: fuente $\to$ base & Acceso & Código \\
    \midrule
    Fed de Nueva York \citep{nyfed2026} & SOFR diario (API \texttt{search.json} con dos fechas) & \% $\to$ decimal & público & \codigo{desk/api/nyfed.py}{fetch_sofr} \\
    Banxico, SIE \citep{banxico2026a,banxico2021} & TIIE de Fondeo a un día, serie \texttt{SF331451} & \% $\to$ decimal & token & \codigo{desk/api/banxico.py}{fetch_series} \\
    Banxico, SIE \citep{banxico2026b} & Tipo de cambio FIX, serie \texttt{SF43718} & MXN por USD, sin conversión & token & \codigo{desk/api/banxico.py}{fetch_series} \\
    FRED \citep{fred2026} & DEXMXUS: pesos por dólar al mediodía de Nueva York (CSV de la serie completa) & MXN por USD, sin conversión & público & \codigo{desk/api/fred.py}{fetch_dexmxus} \\
    MarketDataFeed (SOFR de la Fed de NY y rendimientos par de \citealt{treasury2026}) & Tablas \texttt{sofr\_rates} y \texttt{treasury\_yields} de una base SQLite & \% $\to$ decimal & archivo local & \codigo{desk/api/marketdata.py}{latest_treasury_curve} \\
    \bottomrule
  \end{tabular}
\end{table}

\paragraph{Banxico.} El cliente es \codigo{desk/api/banxico.py}{fetch_series}: pide
\texttt{.../series/<id>/datos/<desde>/<hasta>} con la cabecera \texttt{Bmx-Token}; \texttt{parse\_series}
descarta los datos faltantes (\texttt{"N/E"}) y las filas ilegibles, y si una fecha se repite gana la última. Los
identificadores de serie (\codigo{desk/api/banxico.py}{TIIE_F_SERIES}, \codigo{desk/api/banxico.py}{FIX_SERIES}) se
comprobaron el 2026-11-11 contra el catálogo público del SIE, que no pide token
\citep{banxico2026a,banxico2026b}: SF331451 es «TIIE de Fondeo a Un Día Hábil Bancario, Mediana ponderada por
volumen» y SF43718 es el tipo de cambio pesos por dólar «para solventar obligaciones denominadas en moneda
extranjera, fecha de determinación (FIX)». El token se lee solo de la variable de entorno
\texttt{BANXICO\_TOKEN} (\codigo{desk/api/fixings_fetch.py}{_token}) y nunca entra en una respuesta, un mensaje de
error ni la base.

\paragraph{Fijaciones.} \codigo{desk/api/fixings_fetch.py}{fetch_fixings} es todo-o-nada: descarga el rango
(máximo 366 días, \codigo{desk/api/fixings_fetch.py}{MAX_RANGE_DAYS}), valida cada valor con las reglas de
\codigo{desk/api/fixings.py}{_validate} y solo entonces inserta; un fallo de red o un valor inválido responde
502 sin escribir nada, un token ausente responde 409 y uno rechazado (400, 401 o 403 de Banxico) responde 502 con
\texttt{code: banxico\_token\_rejected}. Responde \texttt{fetched}, \texttt{inserted}, \texttt{unchanged} y
\texttt{conflicts}. \codigo{desk/api/fixings_fetch.py}{list_sources} informa para cada índice su fuente, si está
disponible y su última fecha. Además de la fuente hay alta manual (\codigo{desk/api/fixings.py}{add_fixing}, 409
si el valor difiere del guardado) y carga por CSV con encabezado exacto \texttt{index,date,value}
(\codigo{desk/api/fixings.py}{import_csv}, atómica: un archivo malo no guarda nada).

\paragraph{Mercado vivo.} \codigo{desk/api/refresh.py}{run_refresh} (ruta \texttt{POST /api/market/refresh},
\codigo{desk/api/routes/market.py}{refresh_market}) actualiza \texttt{MKT-LIVE} bajo un candado de proceso (un
solo refresco a la vez; el segundo recibe 409): ejecuta sin shell
\texttt{fetch\_sofr --latest --treasury --db <MDF\_DB>} con tiempo máximo de 120~s
(\codigo{desk/api/refresh.py}{_run_fetch}), descarga DEXMXUS y reconstruye el snapshot con
\codigo{desk/api/live_market.py}{rebuild_live}. MarketDataFeed es un proyecto aparte del repositorio; el
ejecutable obtiene SOFR de \texttt{markets.newyorkfed.org} y los rendimientos de \texttt{home.treasury.gov}.
\codigo{desk/api/marketdata.py}{latest_treasury_curve} abre su base en solo lectura (\texttt{mode=ro}) y exige la
curva completa de la fecha más reciente (no recurre a una fecha anterior si falta un plazo); toda anomalía es
409. \codigo{desk/api/live_market.py}{build_live_data} parte de \texttt{demo\_data()} y reemplaza solo las
cotizaciones SOFR, el spot y las volatilidades de SOFR y USDMXN; \codigo{desk/api/live_market.py}{save_live}
guarda la fila con \texttt{INSERT OR REPLACE} después de armar el bootstrap, así que si algo falla antes no se
escribe nada y cualquier fallo posterior de \texttt{run\_refresh} deja \texttt{MKT-LIVE} como estaba (502).
\texttt{MKT-LIVE} es de solo lectura (editarlo es 409); sus componentes quedan así
(\codigo{desk/api/live_market.py}{_sources}): TIIE-F, superficie, Hull--White y correlaciones, sintéticos;
curva SOFR, \texttt{proxy} con la fecha del Tesoro; spot, \texttt{live} con la fecha de la última observación de
FRED; volatilidades de factores, \texttt{mixed}. El \texttt{as\_of} del snapshot es la fecha de la curva del
Tesoro más reciente, que puede diferir de la del SOFR de un día y de la del spot; cada componente trae su propia
fecha. La base de MarketDataFeed y el ejecutable se configuran con \texttt{MDF\_DB} y \texttt{MDF\_FETCH}
(\codigo{desk/api/app.py}{create_app}); si no se da \texttt{MDF\_FETCH}, se toma
\texttt{examples/fetch\_sofr} junto al directorio de la base de \texttt{MDF\_DB}.

\paragraph{Snapshots editables.} Para tener mercado en varias fechas, \codigo{desk/api/market.py}{roll_snapshot}
copia un snapshot con otro \texttt{id} y \texttt{as\_of} (conserva su procedencia y la copia es editable aunque el
origen sea \texttt{MKT-LIVE}), y \codigo{desk/api/market.py}{update_snapshot} cambia tasas o spot de un
snapshot editable validando que el bootstrap sigue funcionando antes de guardar.

\paragraph{Historia de factores.} La tabla \texttt{factor\_history} lleva su fuente por fila:
\texttt{synthetic}, \texttt{csv}, \texttt{live} (SOFR con el proxy del Tesoro a 5 años y USDMXN de FRED) o
\texttt{fixings} (derivada de las fijaciones cargadas)
(\codigo{desk/api/risk/factor_history.py}{SOURCES}, \codigo{desk/api/risk/factor_history.py}{derive}); los cambios
a más de 5 días naturales de distancia se omiten al derivar y uno fuera de cota detiene la derivación con 409.
El capítulo~\ref{cap:riesgo} explica cómo la usa el VaR histórico.

\subsection{Respaldo}\label{sec:cap13-respaldo}

\codigo{desk/api/backup.py}{backup} copia la base con la API de respaldo en línea de \texttt{sqlite3}
(consistente con el servidor encendido) a \texttt{qfdesk-AAAAMMDD-HHMMSS.sqlite} (hora local), la verifica con
\texttt{PRAGMA integrity\_check} antes de darle su nombre definitivo y con \texttt{--keep N} conserva solo las $N$
copias más recientes entre los archivos cuyo nombre sigue exactamente ese patrón
(\codigo{desk/api/backup.py}{_copies}). \codigo{desk/api/backup.py}{restore} exige el servidor apagado (no lo
detecta) y guarda antes la base actual como \texttt{<base>.antes-de-restaurar-AAAAMMDD-HHMMSS}. Solo respalda
\texttt{desk/data/desk.db}: no cubre la base de MarketDataFeed.

% -------------------------------------------------------------------------------------------------
\section{Ejemplo numérico}\label{sec:cap13-ejemplo}

Los valores siguientes se obtuvieron ejecutando las funciones citadas (no hay script en \texttt{verificacion/}
para este capítulo, que está fuera de los capítulos~2--12); el fragmento final los reproduce.

\paragraph{Tesoro $\to$ OIS.} Con $y_{3\mathrm{M}}=4.00\,\%$ ($D=91\le365$),
$q=0.04\cdot360/365=\num{0.039452}$; con $y_{5\mathrm{Y}}=3.80\,\%$ ($D=1825>365$),
$q=\bigl[(1+0.019)^2-1\bigr]\cdot360/365=0.038361\cdot0.986301=\num{0.037836}$. El pilar de 1 año (365 días) usa
todavía la primera rama; uno de 366 días ya usaría la segunda (con $y=4\,\%$: \num{0.039452} frente a
\num{0.039847}).

\paragraph{Parseo.} Una fila de la Fed de Nueva York con \texttt{percentRate} $=3.95$ se guarda como $0.0395$;
un dato de Banxico \texttt{"7.25"} de la serie SF331451 se guarda como $0.0725$ y uno \texttt{"N/E"} se omite; una
observación de FRED con valor «.» se descarta.

\paragraph{Sonrisa de \texttt{MX-DEMO}.} Con vencimiento de 182 días ($\vol^{\mathrm{atm}}=0.125$) y $K=19.0$:
$k=\ln(19/18.6)=0.021277$ y $\vol=0.125+0.5\cdot0.021277^2+0.03\cdot0.021277=\num{0.125865}$.

\begin{lstlisting}
from desk.api import live_market, nyfed, banxico, fred, seed_market
print(live_market.treasury_to_ois(0.04, 91))     # 0.0394520548
print(live_market.treasury_to_ois(0.038, 1825))  # 0.0378355068
print(nyfed.parse_sofr('{"refRates":[{"effectiveDate":"2026-11-02","type":"SOFR","percentRate":3.95}]}'))
print(banxico.parse_series('{"bmx":{"series":[{"idSerie":"SF331451","datos":'
      '[{"fecha":"02/11/2026","dato":"7.25"},{"fecha":"03/11/2026","dato":"N/E"}]}]}}'))
print(fred.parse_dexmxus("DATE,DEXMXUS\n2026-11-02,18.4\n2026-11-03,.\n"))
print(seed_market.demo_data()["vol_surface"]["vols"][2][4])   # 0.1258646858
\end{lstlisting}

% -------------------------------------------------------------------------------------------------
\section{Validación}\label{sec:cap13-validacion}

Las pruebas del backend simulan el HTTP o leen archivos locales. Parseo y unidades:
\codigo{desk/api/tests/test_nyfed.py}{test_parse_real_fixture_is_ascending_decimal_and_rounded},
\codigo{desk/api/tests/test_banxico.py}{test_parse_tiie_f_to_decimal_ascending_and_skips_missing},
\codigo{desk/api/tests/test_banxico.py}{test_fetch_sends_token_header_exact_url_and_20s_timeout},
\codigo{desk/api/tests/test_fred.py}{test_parse_discards_empty_and_dot_observations},
\codigo{desk/api/tests/test_marketdata.py}{test_latest_treasury_curve_is_complete_and_decimal},
\codigo{desk/api/tests/test_marketdata.py}{test_latest_date_incomplete_does_not_fall_back_to_older_date}.
Conversión y volatilidades de \texttt{MKT-LIVE}:
\codigo{desk/api/tests/test_live_market.py}{test_conversion_reference_values},
\codigo{desk/api/tests/test_live_market.py}{test_conversion_switches_to_semiannual_compounding_after_one_year},
\codigo{desk/api/tests/test_live_market.py}{test_daily_vol_logret_skips_missing_and_joins_valid_neighbours}.
Refresco: \codigo{desk/api/tests/test_refresh.py}{test_failure_is_502_keeps_live_row_and_frees_the_lock}.
Fijaciones: \codigo{desk/api/tests/test_fixings_fetch.py}{test_different_value_is_a_conflict_and_is_not_overwritten},
\codigo{desk/api/tests/test_fixings_fetch.py}{test_repeating_the_fetch_inserts_nothing_and_changes_nothing}.
Procedencia: \codigo{desk/api/tests/test_market_sources.py}{test_data_quality_rule}. Respaldo:
\codigo{desk/api/tests/test_backup.py}{test_respaldo_y_restauracion_conservan_tablas_y_disparadores}.

\paragraph{Comprobaciones contra las fuentes reales (al documentar).} El 2026-11-11 se consultó el catálogo
público del SIE para las series SF331451 y SF43718 y sus títulos coinciden con los que cita el código
\citep{banxico2026a,banxico2026b}; la base local de MarketDataFeed tiene las tablas \texttt{sofr\_rates}
(columnas \texttt{effective\_date}, \texttt{rate}) y \texttt{treasury\_yields} (\texttt{effective\_date},
\texttt{tenor}, \texttt{rate}) con los 13 plazos del Tesoro, de los que el desk lee nueve. No se pudo comprobar
la respuesta de éxito de Banxico con un token real (no hay token): el formato de \texttt{parse\_series} es el
documentado en el código, sin verificar.

% -------------------------------------------------------------------------------------------------
\section{Límites y supuestos}\label{sec:cap13-limites}

\begin{limites}
  \begin{itemize}
    \item \textbf{TIIE-F sintética.} No hay datos libres de TIIE-F: la curva, la volatilidad diaria de TIIE-F, la superficie de volatilidad, los parámetros de Hull--White y las correlaciones son sintéticos en \texttt{MX-DEMO} y también en \texttt{MKT-LIVE}.
    \item \textbf{Curva SOFR aproximada.} Con Treasuries, sin el diferencial entre swaps SOFR y Treasuries (ecuación~\eqref{eq:cap13-tesoro}); el efecto no se midió.
    \item \textbf{Spot.} DEXMXUS de FRED es el tipo de cambio de mediodía de Nueva York, no el FIX de Banxico.
    \item \textbf{Volatilidades.} Hasta 250 cambios diarios, mínimo 60; la de SOFR se mide sobre el rendimiento par de 5 años del Tesoro. Son diarias y por eso la procedencia de las volatilidades de factores sale como mixta.
    \item \textbf{Banxico.} El cliente está listo pero inactivo sin \texttt{BANXICO\_TOKEN} (lo genera el usuario en el portal del SIE, con captcha) y no se probó contra un token real: el formato de la respuesta de éxito no está verificado, y un rango sin observaciones podría responder 404 de Banxico, que hoy se vería como un 502 genérico.
    \item \textbf{Fijaciones.} SOFR son fijaciones reales de la Fed de Nueva York; TIIE-F y FIX solo entran por alta manual, CSV o por Banxico con token. No hay carga automática programada: se piden con el botón o con \texttt{POST /api/fixings/fetch}.
    \item \textbf{Colateral.} \texttt{MKT-LIVE} no trae puntos forward ni base de cross currency swaps, así que con colateral USD se descuenta con TIIE-F y se avisa.
    \item \textbf{Actualización.} Cada «Actualizar datos» reescribe \texttt{MKT-LIVE} por completo; requiere MarketDataFeed compilado y acceso a la red; \texttt{MX-DEMO} no cambia.
    \item \textbf{Fuera de alcance.} Importación de archivos de PiP, Valmer o Bloomberg; curvas reales de TIIE-F; usuarios, roles y aprobación.
    \item \textbf{Respaldo.} Solo de la base del desk; no hay copia fuera de la máquina ni restauración con el servidor encendido.
  \end{itemize}
\end{limites}

% -------------------------------------------------------------------------------------------------
\section{Referencias}\label{sec:cap13-referencias}

Las fuentes de datos son documentación institucional, no artículos: la definición de la TIIE de Fondeo es la nota
técnica de Banxico \citep{banxico2021} y el catálogo de series del SIE \citep{banxico2026a,banxico2026b}; SOFR se
define y se publica en la página de la Fed de Nueva York \citep{nyfed2026}; los rendimientos par son los
publicados por el Tesoro \citep{treasury2026}; el tipo de cambio de mediodía es la serie H.10 de la Reserva
Federal difundida por FRED \citep{fred2026}. Para el papel del FIX y de las fijaciones en el ciclo de vida de una
operación, véase el capítulo~\ref{cap:ciclo-vida}.
\end{sloppypar}
